\section{Dataset Description}\label{sec:dataset-description}
Data from two different datasets called ``Ultra2'' and ``Competition'' are used.
Patients are referred to with the initial letter of their dataset, followed by their ID number, e.g., U17 or C03.

% todo where does this data come from?

Both datasets consist of seizure annotations and \gls{eeg} data, stored as \gls{edf} files.
The seizure's starts, and in some cases the type and end, were contained in the annotations.


% todo
% number of channel, type of channel (distal, proximal)
% explain visits
% which patients were there
% how patients are called

\subsection{Ethics Votum}\label{subsec:ethics-votum}

\subsection{Device Description}\label{subsec:device-description}
% todo die Antragsnummer zu der Studie mit dem UNEEG Device ist: 20-1175

\subsection{Origin of Seizure Annotations}\label{subsec:origin-of-seizure-annotations}

\subsection{Inclusion Criteria}\label{subsec:inclusion-criteria}
For a patient to be included in this study, generally speaking, they needed to have enough \gls{eeg} data and seizures.
These were necessary, so the models have enough examples to learn features which can be used to forecast seizures and for the evaluation to carry statistical weight.
Only patients who had completed the U002 study at the point of writing and thus whose entire \gls{eeg} data was available were considered. % todo correct study name?

Patients needed to have at least 10 valid seizures.
In this context, a seizure is considered valid if no other seizure occurs within its preictal interval (\cref{subsubsec:interval_types}), i.e., a seizure must occur at least $\sim$\SI{65}{\min} after the previous seizure.

Regarding the recording duration, there are two separate concepts.
First, there is the total monitoring timespan, which is the duration from the start of the first recording to the end of the last.
However, data is not recorded all the time, since the device is regularly taken off for a while.
So secondly, the total duration of all available recordings is taking into account.
Based on these, we want to quantify how much of the monitoring timespan data is available.
We define this so-called \textbf{recording coverage} as:

\[\text{recording coverage} = \frac{\text{total recorded duration}}{\text{total monitoring timespan}}\]

Patients with a recording coverage less than \SI{60}{\percent} were excluded.
No patients needed to be excluded on the basis of the total recorded duration, as no patients with little \gls{eeg} data had enough seizures.

\Cref{tab:patient_inclusion} shows which patients were included and why.

\begin{table}[hptb]
    \centering
    \setlength{\tabcolsep}{9pt}
    \renewcommand{\arraystretch}{1.15}
    \rowcolors{2}{gray!10}{white} % Alternating row colors
    \begin{tabular}{|l r|r r r|r r r r|}
\hline
\rowcolor{gray!25} \multicolumn{2}{|c|}{Patient} & \multicolumn{3}{c|}{Seizures} & \multicolumn{4}{c|}{Timespan} \\
\rowcolor{gray!25} ID & Incl. & Total & Valid & Met & Monitoring & Recorded & Coverage & Met \\
\hline
C01 & $\checkmark$ & 92 & 58 & $\checkmark$ & 494 & 388 & \SI{78}{\percent} & $\checkmark$ \\
C02 & $\checkmark$ & 54 & 53 & $\checkmark$ & 537 & 487 & \SI{91}{\percent} & $\checkmark$ \\
C03 & $\checkmark$ & 63 & 63 & $\checkmark$ & 415 & 394 & \SI{95}{\percent} & $\checkmark$ \\
U01 & $\checkmark$ & 259 & 195 & $\checkmark$ & 401 & 381 & \SI{95}{\percent} & $\checkmark$ \\
U02 & $\times$ & 2 & 2 & $\times$ & 293 & 102 & \SI{35}{\percent} & $\times$ \\
U03 & $\checkmark$ & 59 & 52 & $\checkmark$ & 394 & 295 & \SI{75}{\percent} & $\checkmark$ \\
U04 & $\checkmark$ & 64 & 62 & $\checkmark$ & 197 & 180 & \SI{91}{\percent} & $\checkmark$ \\
U05 & $\checkmark$ & 76 & 68 & $\checkmark$ & 358 & 299 & \SI{84}{\percent} & $\checkmark$ \\
U06 & $\times$ & 0 & 0 & $\times$ & 93 & 84 & \SI{90}{\percent} & $\checkmark$ \\
U07 & $\checkmark$ & 13 & 13 & $\checkmark$ & 132 & 124 & \SI{94}{\percent} & $\checkmark$ \\
U08 & $\times$ & 9 & 7 & $\times$ & 179 & 137 & \SI{76}{\percent} & $\checkmark$ \\
U09 & $\times$ & 0 & 0 & $\times$ & 301 & 291 & \SI{97}{\percent} & $\checkmark$ \\
U10 & $\times$ & 0 & 0 & $\times$ & 201 & 165 & \SI{82}{\percent} & $\checkmark$ \\
U11 & $\times$ & 0 & 0 & $\times$ & 216 & 155 & \SI{72}{\percent} & $\checkmark$ \\
U12 & $\checkmark$ & 37 & 30 & $\checkmark$ & 309 & 199 & \SI{64}{\percent} & $\checkmark$ \\
U13 & $\times$ & 4 & 4 & $\times$ & 202 & 115 & \SI{57}{\percent} & $\times$ \\
U14 & $\times$ & 5 & 3 & $\times$ & 420 & 398 & \SI{95}{\percent} & $\checkmark$ \\
U15 & $\checkmark$ & 30 & 23 & $\checkmark$ & 352 & 308 & \SI{88}{\percent} & $\checkmark$ \\
U16 & $\checkmark$ & 16 & 16 & $\checkmark$ & 346 & 222 & \SI{64}{\percent} & $\checkmark$ \\
U17 & $\checkmark$ & 21 & 21 & $\checkmark$ & 264 & 248 & \SI{94}{\percent} & $\checkmark$ \\
U18 & $\times$ & 4 & 4 & $\times$ & 210 & 188 & \SI{89}{\percent} & $\checkmark$ \\
U22 & $\times$ & 3 & 3 & $\times$ & 3 & 3 & \SI{89}{\percent} & $\checkmark$ \\
\hline
\end{tabular}

    \caption[Patient Inclusion]{Patient Inclusion: shows which inclusion criteria were met for which patients. Next to the valid seizures ($\geq 10$) and recording coverage ($\geq\SI{60}{\percent}$) conditions is denoted whether they were met.}
    \label{tab:patient_inclusion}
\end{table}



